\section{Moments}

For an integer $k\ge 0$, the \textbf{$k$-th moment} of a
random variable $X$ is the number $\mathbf{E}(X^k)$.
The \textbf{$k$-th moment around a point $c\in\R$} is the
number $\mathbf{E}(X-c)^k$. If $c$ is not mentioned
it is understood that the moment is around 0. The
\textbf{$k$-th central moment} is the $k$-th moment around
$\mathbf{E} X$ (when it exists), i.e.,
$\mathbf{E}(X-\mathbf{E} X)^k$. In this terminology,
the variance is the second central moment.

The sequence of moments (usually around 0, or around
$\mathbf{E} X$) often contains important information
about the behavior of $X$ and is an important
computational and theoretical tool. Important special
distributions often turn out to have interesting
sequences of moments. Also note that by the
monotonicity of the $L_p$ norms (inequality
\eqref{eq:lpnormmono} in Section \ref{sec-expec-prop}), the
set of values $r\ge 0$ such that $\mathbf{E}(X^r)$
exists (one can also talk about $r$-th moments for
non-integer $r$, but that is much less commonly
discussed) is an interval containing $0$.

A nice characterization of the variance is that it is
the minimal second moment. To compute the second
moment around a point $t$, we can write

\begin{eqnarray*}
\mathbf{E}(X-t)^2 &=& \mathbf{E}[ (X-\mathbf{E} X)-(t-\mathbf{E} X)]^2 \\ &=& \mathbf{E}(X-\mathbf{E} X)^2 + \mathbf{E}(t-\mathbf{E} X)^2 + 2(t-\mathbf{E} X) \mathbf{E}(X-\mathbf{E} X) \\
&=& \mathbf{V}(X) + (t-\mathbf{E} X)^2 \ge \mathbf{V}(X).
\end{eqnarray*}

\noindent So the function $t\to\mathbf{E}(X-t)^2$ is a
quadratic polynomial that attains its minimum at
$t=\mathbf{E} X$, and the value of the minimum is
$\mathbf{V}(X)$. In words, the identity

\[
\mathbf{E}(X-t)^2 = \mathbf{V}(X) + (t-\mathbf{E} X)^2
\]

\noindent says that ``the second moment around $t$ is equal to
the second moment around the mean $\mathbf{E} X$ plus
the square of the distance between $t$ and the
mean''. Note that this is analogous to (and
mathematically equivalent to) the Huygens-Steiner
theorem (also called the Parallel Axis theorem, see
Wikipedia) from mechanics, which says that ``the
moment of inertia of a body with unit mass around a
given axis $L$ is equal to the moment of inertia
around the line parallel to $L$ passing through the
center of mass of the body, plus the square of the
distance between the two lines''. Indeed, the
``moment'' terminology seems to originate in this
physical context.
